\makeatletter
\def\input@path{{preprints-template/preprints_template_latex/}}
\let\preprint@bibliographystyle\bibliographystyle
\renewcommand{\bibliographystyle}[1]{%
  \preprint@bibliographystyle{preprints-template/preprints_template_latex/#1}}
\makeatother

\documentclass[preprints,article,submit,oneauthor]{Definitions/mdpi}

\makeatletter
\let\bibliographystyle\preprint@bibliographystyle
\makeatother

\graphicspath{{preprints-template/preprints_template_latex/}}

\firstpage{1}
\makeatletter
\setcounter{page}{\@firstpage}
\makeatother
\pubvolume{1}
\issuenum{1}
\articlenumber{0}
\pubyear{2026}
\copyrightyear{2026}
\datereceived{}
\daterevised{}
\dateaccepted{}
\datepublished{}

\newcounter{manuscriptresult}[section]
\renewcommand{\themanuscriptresult}{\thesection.\arabic{manuscriptresult}}
\theoremstyle{mdpi}
\newtheorem{theorem}[manuscriptresult]{Theorem}
\newtheorem{lemma}[manuscriptresult]{Lemma}
\newtheorem{proposition}[manuscriptresult]{Proposition}
\newtheorem{corollary}[manuscriptresult]{Corollary}

\theoremstyle{definition}
\newtheorem{remark}[manuscriptresult]{Remark}

\newcommand{\C}{\mathbb C}
\newcommand{\D}{\mathbb D}
\newcommand{\Torus}{\mathbb T}
\newcommand{\Mn}{\mathbb M_n(\C)}
\newcommand{\alg}{\operatorname{alg}}
\newcommand{\distop}{\operatorname{dist}}
\newcommand{\RePart}{\operatorname{Re}}
\newcommand{\diag}{\operatorname{diag}}
\newcommand{\iu}{\mathrm{i}}

\Title{The numerical range is a 2-spectral set}
\Author{Shanmu Jin}
\AuthorNames{Shanmu Jin}
\address{Department of Neurosurgery, Peking Union Medical College
  Hospital, No.~1 Shuaifuyuan, Dongcheng District, Beijing 100730,
  China\\
jinshanmu1129@pku.edu.cn}

\abstract{
For an $n\times n$ complex matrix $A$, let
$W(A)=\{x^*Ax:x\in\C^n,\ \lVert x\rVert_2=1\}$ be its numerical range.
Here $x^*$ denotes the conjugate transpose, $\lVert\,\cdot\,\rVert_2$
the Euclidean norm, and $\lVert\,\cdot\,\rVert$ the induced operator
norm.  For every complex polynomial $p$, we prove that
\[
  \lVert p(A)\rVert\leq
  2\max_{z\in W(A)}\lvert p(z)\rvert
\]
and consequently that $W(A)$ is a $2$-spectral set for $A$.  This
proves Crouzeix's conjecture, and the constant is optimal.  The central
ingredient is a positive-real completion theorem relative to an
auxiliary eigenbasis: an adjoint-algebra constraint on the resolvent
defect of a normalized matrix-valued Carath\'eodory function forces the
underlying matrix to have norm at most $2$.  The classical positive
double-layer calculus supplies such completions for $f(B)$ whenever the
auxiliary matrix $B$ has simple spectrum; the eigenvalues of $f(B)$
may repeat.  Sampling the associated Herglotz kernel at half
the conjugate diagonal entries and at the origin cancels the completion
term and leaves a comparison of two weighted Gramians; a Stein identity
yields the sharp estimate.  Simple-spectrum approximation and convex
outer approximation by domains with continuously differentiable
boundary yield the general theorem.
}

\keyword{Crouzeix conjecture; numerical range; spectral set;
matrix-valued Herglotz kernel; double-layer potential}
\MSC{Primary 47A25; Secondary 47A12, 15A60}

\begin{document}

\section{Introduction}
\label{sec:introduction}

Let $\C$ denote the complex field and, for $n\geq1$, let $\Mn$ denote
the algebra of $n\times n$ complex matrices.  For $A\in\Mn$, its
numerical range is
\begin{equation}
  W(A)=\{x^*Ax:x\in\C^n,\ \lVert x\rVert_2=1\}.
  \label{eq:numerical-range}
\end{equation}
Here $x^*$ denotes the conjugate transpose, $\sigma(A)$ denotes the
spectrum, $\lVert x\rVert_2$ is the Euclidean norm, $\lVert X\rVert$ is
the induced operator norm, and scalar moduli are denoted by
$\lvert\,\cdot\,\rvert$.  The set $W(A)$ is compact, the
Toeplitz--Hausdorff theorem makes it convex, and
$\sigma(A)\subset W(A)$.

A compact set $K\subset\C$ containing $\sigma(A)$ is called a
$C$-spectral set for $A$ if
\begin{equation}
  \lVert r(A)\rVert\leq C\max_{z\in K}\lvert r(z)\rvert
  \label{eq:spectral-set-definition}
\end{equation}
for every rational function $r$ pole-free on $K$.  Our main result
settles Crouzeix's conjecture by identifying the sharp universal
constant for $K=W(A)$.

\begin{theorem}[Main theorem]
\label{thm:main}
Let $n\geq1$, let $A\in\Mn$, and let $p\in\C[z]$.  Then
\begin{equation}
  \lVert p(A)\rVert
  \leq2\max_{z\in W(A)}\lvert p(z)\rvert.
  \label{eq:main-bound}
\end{equation}
Consequently, $W(A)$ is a $2$-spectral set for $A$.
\end{theorem}

The factor $2$ is optimal: for
\begin{equation}
  N=\begin{pmatrix}0&2\\0&0\end{pmatrix}
  \quad\text{and}\quad p(z)=z,
  \label{eq:sharp-example}
\end{equation}
one has $W(N)=\{z\in\C:\lvert z\rvert\leq1\}$ and
$\lVert N\rVert=2$.

For later use, write
\[
  \D=\{z\in\C:\lvert z\rvert<1\},
  \qquad
  \Torus=\partial\D.
\]
We use $I$ for the identity matrix of the size dictated by context,
$X^*$ for the adjoint, put $\RePart X=(X+X^*)/2$, and let $\alg(T)$
denote the unital algebra generated by $T$.  For Hermitian matrices
$X$ and $Y$, we write $X\preceq Y$ (equivalently, $Y\succeq X$) when
$Y-X$ is positive semidefinite.

Crouzeix proved the estimate with constant $2$ for $2\times2$ matrices and
formulated the general conjecture in~\cite{Crouzeix2004}.  His
subsequent dimension-free theorem gave the constant $11.08$~%
\cite{Crouzeix2007}.  After further structural and numerical work~%
\cite{Crouzeix2016}, Crouzeix and Palencia proved that $W(A)$ is a
$(1+\sqrt2)$-spectral set~\cite{CrouzeixPalencia2017}.  Their
double-layer and Cauchy-transform argument was clarified in~%
\cite{RansfordSchwenninger2018} and developed in several directions in~%
\cite{CaldwellGreenbaumLi2018,CrouzeixGreenbaum2019}.

The intervening literature established the conjectured constant for
perturbed Jordan blocks and, more generally, scalar translates of
cyclic weighted shifts, as well as for tridiagonal $3\times3$ matrices
with elliptic numerical range centered at an eigenvalue~%
\cite{ChoiGreenbaum2012,Choi2013,GladerKurulaLindstrom2018}.
Numerical and extremal investigations appear in~%
\cite{GreenbaumOverton2018,Li2020}, and a broad account of the problem
and its variants is given in~\cite{BickelEtAl2020}.  Subsequent work
has treated Blaschke-product level sets, cyclicity reductions,
compressions of shifts, further nilpotent classes, and quantitative
bounds for KMS matrices~%
\cite{BickelGorkin2024,OLoughlinVirtanen2024,BickelEtAl2024,
CrouzeixGreenbaumLi2024}.  An alternative proof for a class of weighted shifts
appears in~\cite{CrouzeixGreenbaum2025}.  Across these cited
reductions, the best dimension-free universal bound remained
$1+\sqrt2$.

\subsection{The sharpness barrier in the symmetrized calculus}

The positive boundary measure used below originates in the
convex-domain integral calculus of Delyon--Delyon~%
\cite{DelyonDelyon1999}; see also~%
\cite{BadeaCrouzeixDelyon2006,CrouzeixPalencia2017}.  If $B$ has
numerical range in a convex domain $\Omega$ with continuously
differentiable boundary and $f$ is bounded by one on $\Omega$, this
calculus produces a unital positive map $\Phi$ and a holomorphic Cauchy
companion $g_f$ such that
\begin{equation}
  2\Phi(f)=f(B)+g_f(B)^*.
  \label{eq:intro-symmetrized}
\end{equation}
Thus positivity directly controls a symmetrized functional calculus,
while the target estimate concerns $f(B)$.  The
Crouzeix--Palencia argument couples the two terms
in~\eqref{eq:intro-symmetrized}; the examples
in~\cite{RansfordSchwenninger2018} show that its abstract norm estimate
is sharp at $1+\sqrt2$.  The corresponding general uniform-algebra
theorem has the same sharp constant; its additional unital case was
proposed as a possible route to $2$
in~\cite{ClouatreOstermannRansford2023}.
Schwenninger and de Vries revisited the double-layer method~%
\cite{SchwenningerDeVries2025}.  Malman et al.\ proved that a
configuration-constant refinement gives a strict improvement for every
fixed numerical-range shape; they also constructed thin quadrilaterals
for which the resulting bounds approach
$1+\sqrt2$~\cite{MalmanEtAl2025}.
Separately,
continuity and compactness give a constant $C_n<1+\sqrt2$ in every
fixed matrix dimension $n$~\cite{MalmanRatio2024}.

The key additional information is algebraic.  For the approximants
constructed in Section~\ref{sec:passage}, the companion
in~\eqref{eq:intro-symmetrized} belongs to the algebra generated by
$B^*$, where the simple-spectrum matrix $B$ supplies an eigenbasis that
also diagonalizes $f(B)$.  Treating $f(B)$ and the companion jointly
preserves this coupling.

\subsection{Auxiliary-basis de-symmetrization}

Our argument uses the whole Cayley family associated
with~\eqref{eq:intro-symmetrized}.  It has four steps.

\begin{enumerate}
\item Applying $\Phi$ to
\[
  \zeta\longmapsto\frac{1+w f(\zeta)}{1-w f(\zeta)},
  \qquad w\in\D,
\]
produces a matrix-valued Carath\'eodory function $H$, meaning that
$H(0)=I$ and $\RePart H(w)\succeq0$.

\item For a simple-spectrum auxiliary matrix $B$, put $T=f(B)$.  The
double-layer identity gives the positive-real completion
\begin{equation}
  H(w)-(I-wT)^{-1}\in\alg(B^*),\qquad w\in\D.
  \label{eq:intro-completion}
\end{equation}

\item The eigenbasis of $B$ diagonalizes $T=f(B)$ even when some values
of $f$ on the spectrum of $B$ coincide, and it turns the unknown term
in~\eqref{eq:intro-completion} into a diagonal analytic correction.  We
sample the Herglotz kernel of $H$ at half the conjugate diagonal entries
of $T$ and add one vector sample at the origin.  Repeated sampling
points are allowed, and the origin sample cancels the diagonal
correction exactly.

\item The surviving positive-matrix inequality compares two weighted
Gramians.  Its difference is positive, and an eigenvector argument
followed by a Stein identity yields $\lVert T\rVert\leq2$.
\end{enumerate}

The proof uses two limits: simple-spectrum matrices tend to $A$,
followed by outer domains converging to $W(A)$.  Together they yield
the general case.

\section{Positive-real kernels}
\label{sec:herglotz}

Let $\mathcal K$ be an $\Mn$-valued kernel on $\D\times\D$.  We call
$\mathcal K$ positive if, for every finite choice
$w_0,\ldots,w_m\in\D$ and $\xi_0,\ldots,\xi_m\in\C^n$,
\begin{equation}
  \sum_{i,j=0}^m
  \xi_i^*\mathcal K(w_i,w_j)\xi_j\geq0.
  \label{eq:positive-kernel-definition}
\end{equation}

We shall use the following standard matrix-valued form of the Herglotz
theorem.  Its kernel consequence is included to fix the normalization.

\begin{lemma}[Matrix Herglotz kernel]
\label{lem:herglotz-kernel}
Let $F\colon\D\to\Mn$ be analytic with
$\RePart F(w)\succeq0$ for every $w\in\D$.  Then
\begin{equation}
  \mathcal L_F(w,z)
  =\frac{F(w)+F(z)^*}{1-w\overline z}
  \label{eq:herglotz-kernel}
\end{equation}
is a positive kernel.
\end{lemma}

\begin{proof}
The matrix Herglotz representation gives a positive
$\Mn$-valued measure $M$ on $\Torus$, and a Hermitian matrix $J_F$, such
that
\begin{equation}
  F(w)=\iu J_F+\int_{\Torus}\frac{\zeta+w}{\zeta-w}\,dM(\zeta).
  \label{eq:matrix-herglotz-representation}
\end{equation}
Indeed, for each $u\in\C^n$ the scalar Herglotz theorem gives a positive
finite measure $\mu_u$ on $\Torus$ and a real number $c_u$ such that
\[
  u^*F(w)u=\iu c_u+
  \int_{\Torus}\frac{\zeta+w}{\zeta-w}\,d\mu_u(\zeta).
\]
Uniqueness in the scalar theorem shows that, for every Borel set
$E\subset\Torus$, the map $u\mapsto\mu_u(E)$ is a nonnegative quadratic
form.  Indeed, uniqueness applied to the corresponding identities for
the scalar functions gives, as identities of measures,
\[
  \mu_{\alpha u}=\lvert\alpha\rvert^2\mu_u,
  \qquad
  \mu_{u+v}+\mu_{u-v}=2\mu_u+2\mu_v
  \quad(\alpha\in\C,\ u,v\in\C^n).
\]
Its polarization defines
\[
  x^*M(E)y=\frac14\sum_{k=0}^3 \iu^{-k}
  \mu_{x+\iu^k y}(E)\qquad(x,y\in\C^n).
\]
Thus $M$ is a finite $\Mn$-valued measure and
$u^*M(E)u=\mu_u(E)\geq0$, so $M$ is positive.  The scalar imaginary
constants polarize to the Hermitian matrix
$J_F=(F(0)-F(0)^*)/(2\iu)$, and the scalar representations assemble to~%
\eqref{eq:matrix-herglotz-representation}.

For $\lvert\zeta\rvert=1$, direct calculation gives
\begin{equation}
  \frac{1}{1-w\overline z}\left(
   \frac{\zeta+w}{\zeta-w}
   +\overline{\frac{\zeta+z}{\zeta-z}}
  \right)
  =\frac{2}{(1-w\overline\zeta)(1-\overline z\zeta)}.
  \label{eq:scalar-herglotz-factorization}
\end{equation}
The constant $\iu J_F$ cancels from~\eqref{eq:herglotz-kernel}.  Substituting~%
\eqref{eq:matrix-herglotz-representation} and then summing as in~%
\eqref{eq:positive-kernel-definition} yields
\[
  2\int_{\Torus}
  \left(\sum_i\frac{\xi_i}{1-\overline{w_i}\zeta}\right)^*
  dM(\zeta)
  \left(\sum_j\frac{\xi_j}{1-\overline{w_j}\zeta}\right)
  \geq0.\qedhere
\]
\end{proof}

\section{The positive-real completion principle}
\label{sec:completion}

The following completion theorem is the core of the argument: it
converts a positive-real completion constrained by the adjoint algebra
of an auxiliary simple-spectrum matrix into the sharp norm bound.

\begin{theorem}[Positive-real completion]
\label{thm:completion}
Let
\begin{equation}
  B=S\diag(\beta_1,\ldots,\beta_n)S^{-1},
  \qquad
  T=S\Lambda S^{-1},
  \qquad
  \Lambda=\diag(\lambda_1,\ldots,\lambda_n),
  \label{eq:T-diagonalization}
\end{equation}
where $S$ is invertible, the numbers $\beta_1,\ldots,\beta_n$ are
distinct, and $\lvert\lambda_i\rvert\leq1$ for every $i$.  The numbers
$\lambda_i$ may repeat.
Let $H\colon\D\to\Mn$ be analytic and satisfy
\begin{equation}
  H(0)=I,\qquad \RePart H(w)\succeq0\quad(w\in\D),
  \label{eq:completion-positive-real}
\end{equation}
and
\begin{equation}
  H(w)-(I-wT)^{-1}\in\alg(B^*)\qquad(w\in\D).
  \label{eq:completion-algebra-condition}
\end{equation}
Then $\lVert T\rVert\leq2$.
\end{theorem}

\begin{proof}
Use the simultaneous diagonalizations in~\eqref{eq:T-diagonalization}
and put $G=S^*S$.  Since the $\beta_i$ are distinct, polynomial
interpolation shows that
\begin{equation}
  \alg(B^*)
  =\left\{(S^{-1})^* D S^* \,:\,
  D\ \text{is diagonal}\right\}.
  \label{eq:adjoint-algebra-diagonal}
\end{equation}
Set
\[
  \Delta(w)=H(w)-(I-wT)^{-1}.
\]
By~\eqref{eq:adjoint-algebra-diagonal}, for each $w\in\D$ there is a
unique diagonal matrix $\Theta(w)$ such that
\[
  \Delta(w)=(S^{-1})^*\Theta(w)S^*.
\]
In fact,
\[
  \Theta(w)=S^*\Delta(w)(S^*)^{-1},
\]
so $\Theta$ is analytic; moreover, $\Theta(0)=0$ because $H(0)=I$.
Write $\Theta(w)=\diag(\theta_1(w),\ldots,\theta_n(w))$.  On setting
$\widetilde H(w)=S^*H(w)S$, we obtain
\begin{equation}
  \widetilde H(w)=G(I-w\Lambda)^{-1}+\Theta(w)G.
  \label{eq:Htilde-decomposition}
\end{equation}
Since $\RePart\widetilde H(w)=S^*(\RePart H(w))S$, the real part of
$\widetilde H(w)$ is positive semidefinite.  Consequently,
Lemma~\ref{lem:herglotz-kernel} applies to $\widetilde H$.

Define Hermitian matrices $P,Q,Y\in\Mn$ by
\begin{equation}
  P_{ij}=\frac{G_{ij}}{1-\overline{\lambda_i}\lambda_j/4},
  \qquad
  Q_{ij}=\frac{G_{ij}}{1-\overline{\lambda_i}\lambda_j/2},
  \qquad
  Y=Q-P.
  \label{eq:PQY-definition}
\end{equation}
Let $e_1,\ldots,e_n$ be the standard basis of $\C^n$, and fix
$u=(u_1,\ldots,u_n)^T\in\C^n$.  In the positive-kernel inequality use
the $n$ points and vectors
\begin{equation}
  w_i=\frac{\overline{\lambda_i}}2,\qquad
  \xi_i=u_i e_i\qquad(1\leq i\leq n),
  \label{eq:eigenvalue-samples}
\end{equation}
and add the point $w_0=0$ with vector
\begin{equation}
  \xi_0=v=-G^{-1}Pu.
  \label{eq:compensating-vector}
\end{equation}
The points $w_i$ may coincide.  Kernel positivity applies to arbitrary
finite indexed lists of points and vectors, including $w_i=w_0$ when
$\lambda_i=0$.
We first verify that this sample cancels the entire unknown function
$\Theta$.  The contribution of the second term in~%
\eqref{eq:Htilde-decomposition} to the kernel quadratic form is
\begin{equation}
  2\RePart\sum_{i=1}^n
  \overline{u_i}\theta_i(w_i)
  \left(
    (Gv)_i+
    \sum_{j=1}^n
    \frac{G_{ij}u_j}{1-w_i\overline{w_j}}
  \right).
  \label{eq:theta-contribution}
\end{equation}
Because
\[
  1-w_i\overline{w_j}
  =1-\overline{\lambda_i}\lambda_j/4,
\]
the parenthesis in~\eqref{eq:theta-contribution} is
$(Gv+Pu)_i$, which vanishes by~%
\eqref{eq:compensating-vector}.

It remains to calculate the contribution of
$G(I-w\Lambda)^{-1}$.  Between $w_i$ and $w_j$ its $(i,j)$ entry is
\begin{equation}
  \frac{2G_{ij}}
  {(1-\overline{\lambda_i}\lambda_j/2)
   (1-\overline{\lambda_i}\lambda_j/4)}
  =(4Q-2P)_{ij}.
  \label{eq:sampled-main-block}
\end{equation}
Between the eigenvalue samples and the origin, the two blocks are
$G+Q$, and the block at the origin is $2G$.  After the cancellation of
$\Theta$, kernel positivity therefore gives
\begin{equation}
\begin{aligned}
0&\leq
\begin{pmatrix}u\\v\end{pmatrix}^{\!*}
\begin{pmatrix}
  4Q-2P&G+Q\\
  G+Q&2G
\end{pmatrix}
\begin{pmatrix}u\\v\end{pmatrix}\\
&=u^*\bigl((4Q-2P)-(G+Q)G^{-1}P
-PG^{-1}(G+Q)+2PG^{-1}P\bigr)u\\
&=u^*\bigl(4Y-YG^{-1}P-PG^{-1}Y\bigr)u,
\end{aligned}
\label{eq:pre-gramian-inequality}
\end{equation}
where in the second line we used $v=-G^{-1}Pu$, and in the last line
we used $Q=P+Y$.
Since $u$ was arbitrary,
\begin{equation}
  4Y-YG^{-1}P-PG^{-1}Y\succeq0.
  \label{eq:Y-inequality}
\end{equation}

Balance the diagonalization by defining
\begin{equation}
\begin{aligned}
  \widetilde T&=G^{1/2}\Lambda G^{-1/2},
  &\widehat P&=G^{-1/2}PG^{-1/2},\\
  \widehat Q&=G^{-1/2}QG^{-1/2},
  &\widehat Y&=\widehat Q-\widehat P.
\end{aligned}
\label{eq:balanced-definitions}
\end{equation}
Expanding the scalar kernels in~\eqref{eq:PQY-definition} into
geometric series gives
\begin{equation}
  \widehat P=\sum_{k=0}^{\infty}
  4^{-k}\widetilde T^{*k}\widetilde T^k,
  \qquad
  \widehat Q=\sum_{k=0}^{\infty}
  2^{-k}\widetilde T^{*k}\widetilde T^k.
  \label{eq:two-gramians}
\end{equation}
Indeed, $\widetilde T^k=G^{1/2}\Lambda^k G^{-1/2}$, so the sequence
$(\widetilde T^k)$ is bounded even when some
$\lvert\lambda_j\rvert=1$; the
geometric weights give norm convergence.  It follows that
\begin{equation}
  \widehat Y=
  \sum_{k=1}^{\infty}(2^{-k}-4^{-k})
  \widetilde T^{*k}\widetilde T^k\succeq0.
  \label{eq:Yhat-positive}
\end{equation}
Taking the congruence of~\eqref{eq:Y-inequality} by $G^{-1/2}$ yields
\begin{equation}
  4\widehat Y-\widehat Y\widehat P
  -\widehat P\widehat Y\succeq0.
  \label{eq:Yhat-inequality}
\end{equation}

Let $e$ be an eigenvector of the Hermitian matrix $\widehat P$, say
$\widehat Pe=\alpha e$.  If $\alpha>2$, then~%
\eqref{eq:Yhat-inequality} gives
\begin{equation}
  0\leq2(2-\alpha)e^*\widehat Ye.
  \label{eq:eigenvector-evaluation}
\end{equation}
Since $\widehat Y\succeq0$ and $\alpha>2$,~%
\eqref{eq:eigenvector-evaluation} forces $e^*\widehat Ye=0$.  Using the
series in~\eqref{eq:Yhat-positive}, we obtain
\[
  0=e^*\widehat Ye
  =\sum_{k=1}^{\infty}(2^{-k}-4^{-k})
    \lVert\widetilde T^k e\rVert_2^2
  \geq\frac14\lVert\widetilde Te\rVert_2^2.
\]
Thus $\widetilde Te=0$, and the first series in~%
\eqref{eq:two-gramians} gives $\widehat Pe=e$, contradicting
$\widehat Pe=\alpha e$ with $\alpha>2$.  The first series in~%
\eqref{eq:two-gramians} also begins with $I$, and hence
\begin{equation}
  I\preceq\widehat P\preceq2I.
  \label{eq:Phat-bounds}
\end{equation}

Finally, the first Gramian satisfies the Stein identity
\begin{equation}
  \widehat P
  =I+\frac14\widetilde T^*\widehat P\widetilde T.
  \label{eq:stein-identity}
\end{equation}
Using~\eqref{eq:Phat-bounds}, we obtain
\[
  \widetilde T^*\widetilde T
  \preceq \widetilde T^*\widehat P\widetilde T
  =4(\widehat P-I)\preceq4I.
\]
Thus $\lVert\widetilde T\rVert\leq2$.  The polar decomposition
$S=UG^{1/2}$ gives
\[
  T=U\widetilde TU^*,
\]
where $U$ is unitary.  Therefore $\lVert T\rVert\leq2$.
\end{proof}

\begin{remark}
\label{rem:completion-novelty}
The factor $1/2$ in the sampling points~%
\eqref{eq:eigenvalue-samples} simultaneously produces the weights
$1/2$ and $1/4$ in~\eqref{eq:two-gramians}.  Their positive difference
and the Stein identity are what recover the sharp constant.  The
origin sample cancels the unknown analytic correction by an exact
algebraic identity.  Distinct auxiliary eigenvalues identify
$\alg(B^*)$ with the diagonal algebra in the chosen basis, while the
diagonal entries of $T$ may repeat.
\end{remark}

\section{The double-layer completion}
\label{sec:double-layer}

We now construct from the numerical range the function required by
Theorem~\ref{thm:completion}.  The measure and companion transform used
here are classical~%
\cite{DelyonDelyon1999,Crouzeix2007,CrouzeixPalencia2017}.  The full
Cayley family yields the algebraic
constraint~\eqref{eq:completion-algebra-condition} relative to the
auxiliary matrix.

\begin{lemma}[Double-layer positive map]
\label{lem:double-layer-map}
Let $B\in\Mn$, and let $\Omega\subset\C$ be a bounded convex domain
with continuously differentiable boundary and
$W(B)\subset\Omega$.  Writing $C(\partial\Omega)$ for the continuous
complex-valued functions on $\partial\Omega$, there is a unital positive map
\[
  \Phi\colon C(\partial\Omega)\longrightarrow\Mn
\]
such that, whenever $h$ is holomorphic on a neighborhood of
$\overline\Omega$,
\begin{equation}
  2\Phi(h)=h(B)+g_h(B)^*,
  \label{eq:double-layer-identity}
\end{equation}
where
\begin{equation}
  g_h(z)=\frac{1}{2\pi\iu}
  \int_{\partial\Omega}
  \frac{\overline{h(\zeta)}}{\zeta-z}\,d\zeta,
  \qquad z\in\Omega.
  \label{eq:companion-transform}
\end{equation}
\end{lemma}

\begin{proof}
Orient $\partial\Omega$ counterclockwise.  Let $\nu(\zeta)$ be the
outward unit normal, let $ds$ be arclength, and put
\[
  Z_\zeta=\zeta I-B,
  \qquad
  \mathcal R_\zeta=Z_\zeta^{-1}.
\]
This resolvent exists because
$\sigma(B)\subset W(B)\subset\Omega$.  Define the matrix-valued
boundary measure
\begin{equation}
  d\mu_B(\zeta)=\frac{1}{2\pi}
  \left(
    \nu(\zeta)\mathcal R_\zeta+
    \overline{\nu(\zeta)}\mathcal R_\zeta^*
  \right)ds.
  \label{eq:double-layer-measure}
\end{equation}
It is positive.  In fact, the supporting-line property of the convex
set $\overline\Omega$ gives
\begin{equation}
  \nu(\zeta)Z_\zeta^*+\overline{\nu(\zeta)}Z_\zeta
  =2\RePart\bigl(\overline{\nu(\zeta)}Z_\zeta\bigr)\succeq0.
  \label{eq:supporting-half-plane}
\end{equation}
Indeed, for every unit vector $x\in\C^n$, the quadratic form on the
left is
\[
  2\RePart\bigl(\overline{\nu(\zeta)}
  (\zeta-x^*Bx)\bigr)\geq0,
\]
because $x^*Bx\in W(B)\subset\Omega$.
Multiplication on the left by $\mathcal R_\zeta^*$ and on the right by
$\mathcal R_\zeta$ turns the left side of~%
\eqref{eq:supporting-half-plane} into the matrix in parentheses in~%
\eqref{eq:double-layer-measure}.

Along the counterclockwise boundary,
$d\zeta=\iu\nu(\zeta)ds$.  The matrix Cauchy formula gives
\begin{equation}
  \frac{1}{2\pi}\int_{\partial\Omega}
  \nu(\zeta)\mathcal R_\zeta\,ds=I.
  \label{eq:first-layer-mass}
\end{equation}
Taking adjoints gives the same identity for the second layer.  Hence
\begin{equation}
  \int_{\partial\Omega}d\mu_B=2I.
  \label{eq:double-layer-mass}
\end{equation}
For $\varphi\in C(\partial\Omega)$, define
\begin{equation}
  \Phi(\varphi)=\frac12
  \int_{\partial\Omega}\varphi(\zeta)\,d\mu_B(\zeta).
  \label{eq:Phi-definition}
\end{equation}
By~\eqref{eq:double-layer-mass}, $\Phi$ is unital and positive, hence
$*$-preserving; moreover,
$\lVert\Phi(\varphi)\rVert\leq\lVert\varphi\rVert_\infty$, where
\[
  \lVert\varphi\rVert_\infty
  =\max_{\zeta\in\partial\Omega}\lvert\varphi(\zeta)\rvert.
\]

For holomorphic $h$, the first layer of $2\Phi(h)$ is $h(B)$ by
Cauchy's formula.  The function $g_h$ in~%
\eqref{eq:companion-transform} is holomorphic on $\Omega$, so
$g_h(B)$ is defined by the holomorphic functional calculus.  To verify
the relevant integral formula using only the holomorphy of $g_h$ on
$\Omega$, choose a positively oriented contour
$\Gamma\subset\Omega$ enclosing $\sigma(B)$.  Cauchy's formula and
Fubini's theorem give
\begin{equation}
\begin{split}
 g_h(B)
 &=\frac{1}{2\pi\iu}\int_{\Gamma}
   g_h(z)(zI-B)^{-1}\,dz\\
 &=\frac{1}{2\pi\iu}\int_{\partial\Omega}
   \overline{h(\zeta)}(\zeta I-B)^{-1}\,d\zeta.
\end{split}
\label{eq:companion-functional-calculus}
\end{equation}
Taking adjoints in~\eqref{eq:companion-functional-calculus} and using
$d\overline\zeta=-\iu\overline{\nu(\zeta)}ds$ gives the second layer of
$2\Phi(h)$.  This proves~\eqref{eq:double-layer-identity}.
\end{proof}

\begin{proposition}[Cayley completion for an auxiliary pair]
\label{prop:cayley-completion}
Let $B$ and $\Omega$ be as in Lemma~\ref{lem:double-layer-map}, and let $f$ be
holomorphic on a neighborhood of $\overline\Omega$ and satisfy
\begin{equation}
  \max_{\zeta\in\overline\Omega}\lvert f(\zeta)\rvert\leq1.
  \label{eq:f-bounded-by-one}
\end{equation}
Put $T=f(B)$.  Then there is an analytic $H\colon\D\to\Mn$ satisfying
\begin{equation}
  H(0)=I,\qquad \RePart H(w)\succeq0\quad(w\in\D),
  \label{eq:cayley-H-positive}
\end{equation}
and
\begin{equation}
  H(w)-(I-wT)^{-1}\in\alg(B^*)\qquad(w\in\D).
  \label{eq:cayley-H-completion}
\end{equation}
Moreover, $\sigma(T)\subset\overline\D$.
\end{proposition}

\begin{proof}
By~\eqref{eq:f-bounded-by-one} and spectral mapping from
$\sigma(B)\subset\Omega$, we have $\sigma(T)\subset\overline\D$.
For $w\in\D$, define on the boundary
\begin{equation}
  c_w(\zeta)=\frac{1+w f(\zeta)}{1-w f(\zeta)}
  \label{eq:cayley-boundary-function}
\end{equation}
and put
\begin{equation}
  H(w)=\Phi(c_w).
  \label{eq:H-from-Phi}
\end{equation}
The map $w\mapsto c_w$ is analytic with values in
$C(\partial\Omega)$, locally uniformly on $\D$.  Hence $H$ is
analytic, and $H(0)=I$.  Furthermore,
\[
  \RePart c_w(\zeta)
  =\frac{1-\lvert w f(\zeta)\rvert^2}
  {\lvert1-w f(\zeta)\rvert^2}\geq0.
\]
The positivity and $*$-preservation of $\Phi$ give
$\RePart H(w)=\Phi(\RePart c_w)\succeq0$.

The boundary expansion, convergent in $C(\partial\Omega)$ locally
uniformly for $w\in\D$,
\begin{equation}
  c_w=1+2\sum_{m=1}^{\infty}w^m f^m
  \label{eq:cayley-series}
\end{equation}
and~\eqref{eq:double-layer-identity} yield, with norm convergence,
\begin{equation}
  H(w)=I+\sum_{m=1}^{\infty}w^m
  \bigl(T^m+g_{f^m}(B)^*\bigr).
  \label{eq:H-series}
\end{equation}
Since the spectral radius of $wT$ is less than $1$, the Neumann series
for $(I-wT)^{-1}$ converges in norm.  Subtracting it from~%
\eqref{eq:H-series} gives
\begin{equation}
  H(w)-(I-wT)^{-1}
  =\sum_{m=1}^{\infty}w^m g_{f^m}(B)^*.
  \label{eq:completion-series}
\end{equation}
Each $g_{f^m}$ is holomorphic on a neighborhood of $\sigma(B)$.
Finite-dimensional holomorphic functional calculus, equivalently
Hermite interpolation followed by Cayley--Hamilton, gives
$g_{f^m}(B)\in\alg(B)$.  Hence
\[
  g_{f^m}(B)^*\in\alg(B^*).
\]
Every partial sum on the right of~\eqref{eq:completion-series}
therefore belongs to $\alg(B^*)$.  This algebra is finite-dimensional
and hence norm closed, which proves~%
\eqref{eq:cayley-H-completion}.
\end{proof}

\begin{corollary}[Auxiliary sharp bound]
\label{cor:auxiliary-bound}
In Proposition~\ref{prop:cayley-completion}, if $B$ has simple
spectrum, then $\lVert f(B)\rVert\leq2$.
\end{corollary}

\begin{proof}
Write
\[
  B=S\diag(\beta_1,\ldots,\beta_n)S^{-1}.
\]
Then
\[
  T=f(B)
  =S\diag\bigl(f(\beta_1),\ldots,f(\beta_n)\bigr)S^{-1}.
\]
Since $\sigma(B)\subset\Omega$, the bound on $f$ gives
$\lvert f(\beta_i)\rvert\leq1$ for every $i$.
Apply Theorem~\ref{thm:completion} to these simultaneous
diagonalizations and to the function $H$ constructed in
Proposition~\ref{prop:cayley-completion}.  The values $f(\beta_i)$ may
repeat.
\end{proof}

\section{Passage to arbitrary matrices and proof of the main theorem}
\label{sec:passage}

Successive approximations pass from the auxiliary setting to arbitrary
matrices.

\begin{lemma}[Simple spectrum and numerical range]
\label{lem:simple-spectrum-density}
Matrices with $n$ distinct eigenvalues are dense in $\Mn$.  Moreover,
for $A,B\in\Mn$,
\begin{equation}
  W(B)\subset
  W(A)+\{z\in\C:\lvert z\rvert\leq\lVert B-A\rVert\}.
  \label{eq:numerical-range-Lipschitz}
\end{equation}
\end{lemma}

\begin{proof}
The discriminant of the characteristic polynomial is a nonzero
polynomial in the matrix entries, as is seen at
$\diag(1,\ldots,n)$.  Its nonvanishing set is therefore dense in
$\Mn$, proving the first assertion.  For a unit vector $x$,
\[
  \lvert x^*(B-A)x\rvert\leq\lVert B-A\rVert,
\]
which proves~\eqref{eq:numerical-range-Lipschitz}.
\end{proof}

\begin{lemma}[Outer approximation]
\label{lem:smooth-outer-approximation}
If $K\subset\C$ is nonempty, compact, and convex, then there are
bounded convex domains $\Omega_j$ with continuously differentiable
boundaries such that
\begin{equation}
  K\subset\Omega_j,
  \qquad
  d_{\mathrm H}(\overline\Omega_j,K)\longrightarrow0,
  \label{eq:outer-Hausdorff}
\end{equation}
where $d_{\mathrm H}$ denotes the Hausdorff distance between compact
sets.
\end{lemma}

\begin{proof}
For $\delta>0$, let
\[
  \Omega_\delta=\{z\in\C:\distop(z,K)<\delta\}.
\]
This is a bounded convex domain and its closure is the Minkowski sum
$K+\delta\overline\D$, so its Hausdorff distance from $K$ is $\delta$.
Identify $\C$ with $\mathbb R^2$.  The metric projection
$\pi_K\colon\C\to K$ is unique and continuous.  Put
$d_K(z)=\distop(z,K)$.  For $z,h\in\C$, comparison first with $\pi_K(z)$
and then with $\pi_K(z+h)$ gives
\[
\begin{split}
  2\RePart\bigl(\overline{z+h-\pi_K(z+h)}\,h\bigr)-\lvert h\rvert^2
  &\leq d_K(z+h)^2-d_K(z)^2\\
  &\leq
  2\RePart\bigl(\overline{z-\pi_K(z)}\,h\bigr)+\lvert h\rvert^2.
\end{split}
\]
The continuity of $\pi_K$ now shows that $d_K^2$ is continuously
differentiable, with real gradient $2(z-\pi_K(z))$.  On the level set
$d_K(z)=\delta$ this gradient has norm $2\delta>0$.  Since
$\partial\Omega_\delta=\{d_K=\delta\}$, the
implicit-function theorem shows that $\partial\Omega_\delta$ is
continuously differentiable.  Taking
$\Omega_j=\Omega_{1/j}$ for $j\geq1$ proves the lemma.
\end{proof}

\begin{proof}[Proof of Theorem~\ref{thm:main}]
Let $A$ and $p$ be arbitrary as in Theorem~\ref{thm:main}.  Choose a
bounded convex domain $\Omega$ with continuously differentiable boundary
such that
\begin{equation}
  W(A)\subset\Omega.
  \label{eq:WA-in-Omega}
\end{equation}
We first prove
\begin{equation}
  \lVert p(A)\rVert
  \leq2\max_{z\in\overline\Omega}\lvert p(z)\rvert.
  \label{eq:fixed-Omega-bound}
\end{equation}
Because $W(A)$ is compactly contained in $\Omega$,
\[
  \distop(W(A),\C\setminus\Omega)>0.
\]
Lemma~\ref{lem:simple-spectrum-density} and~%
\eqref{eq:numerical-range-Lipschitz} therefore give a sequence of
matrices $B_k\to A$ having simple spectrum and satisfying
$W(B_k)\subset\Omega$.  Since $p(B_k)\to p(A)$, it is enough to prove~%
\eqref{eq:fixed-Omega-bound} for an arbitrary such matrix $B$.

Set
\begin{equation}
  m_\Omega=\max_{z\in\overline\Omega}\lvert p(z)\rvert.
  \label{eq:m-Omega}
\end{equation}
If $m_\Omega=0$, the polynomial vanishes on the open set $\Omega$ and
is identically zero.  When $m_\Omega>0$, put
\begin{equation}
  f=\frac{p}{m_\Omega}.
  \label{eq:normalized-f}
\end{equation}
Then
$\max_{z\in\overline\Omega}\lvert f(z)\rvert\leq1$.
Corollary~\ref{cor:auxiliary-bound}, applied directly to $B$ and $f$,
gives
\[
  \lVert f(B)\rVert\leq2.
\]
Equivalently,
\begin{equation}
  \lVert p(B)\rVert\leq2m_\Omega.
  \label{eq:B-Omega-bound}
\end{equation}
Letting the matrices $B_k$ with simple spectrum tend to $A$ proves~%
\eqref{eq:fixed-Omega-bound}.

Apply Lemma~\ref{lem:smooth-outer-approximation} to the compact convex
set $K=W(A)$, obtaining $\Omega_j$.  All these closures eventually
lie in one fixed compact neighborhood of $K$.  Uniform continuity of
$p$ there and~\eqref{eq:outer-Hausdorff} give
\begin{equation}
  \max_{z\in\overline\Omega_j}\lvert p(z)\rvert
  \longrightarrow
  \max_{z\in W(A)}\lvert p(z)\rvert.
  \label{eq:maxima-converge}
\end{equation}
Apply~\eqref{eq:fixed-Omega-bound} to each $\Omega_j$ and pass to the
limit.  This proves~\eqref{eq:main-bound}.

It remains to justify the spectral-set conclusion.  Let $r$ be
rational and pole-free on $W(A)$.  Since $W(A)$ is compact, choose
$\delta>0$ so that $r$ is holomorphic on an open neighborhood of the
closed $\delta$-neighborhood $K_\delta$ of $W(A)$.  Since
$\C\setminus K_\delta$ is connected, polynomial Runge approximation
gives polynomials $p_j$ converging uniformly to $r$ on $K_\delta$.  As
$\sigma(A)\subset W(A)\subset\operatorname{int}K_\delta$, choose a
positively oriented contour $\Gamma\subset\operatorname{int}K_\delta$
enclosing $\sigma(A)$.  Uniform convergence on $\Gamma$, together with
the Cauchy integral formula for the holomorphic functional calculus,
gives $p_j(A)\to r(A)$.  Also
\[
  \max_{z\in W(A)}\lvert p_j(z)\rvert
  \longrightarrow
  \max_{z\in W(A)}\lvert r(z)\rvert.
\]
Apply the polynomial estimate to $p_j$ and pass to the limit to obtain~%
\eqref{eq:spectral-set-definition} with $K=W(A)$ and $C=2$.
\end{proof}

\section{Consequences and scope}
\label{sec:consequences}

Theorem~\ref{thm:main} immediately gives the corresponding estimate for
functions holomorphic on a neighborhood of the numerical range.

\begin{corollary}
\label{cor:holomorphic-bound}
Let $A\in\Mn$, and let $f$ be holomorphic on a neighborhood of
$W(A)$.  Then
\begin{equation}
  \lVert f(A)\rVert
  \leq2\max_{z\in W(A)}\lvert f(z)\rvert.
  \label{eq:holomorphic-bound}
\end{equation}
\end{corollary}

\begin{proof}
Choose $\delta>0$ so that the closed $\delta$-neighborhood of $W(A)$
lies in the open set on which $f$ is holomorphic.  Polynomial Runge
approximation on this compact convex neighborhood, followed by the
Cauchy integral formula for the holomorphic functional calculus, gives
the assertion.
\end{proof}

The Hilbert-space extension follows by the standard finite-dimensional
compression argument; compare~\cite{BickelEtAl2020}.

\begin{corollary}[Hilbert-space form]
\label{cor:operator-bound}
Let $\mathcal H\neq\{0\}$ be a complex Hilbert space, with inner
product linear in the first variable, and let $\mathcal B(\mathcal H)$
denote its algebra of bounded operators.  For
$A\in\mathcal B(\mathcal H)$, put
\[
  W(A)=\{\langle Ax,x\rangle:x\in\mathcal H,\ \lVert x\rVert=1\}.
\]
Then, for every $p\in\C[z]$,
\begin{equation}
  \lVert p(A)\rVert
  \leq 2\sup_{z\in W(A)}\lvert p(z)\rvert.
  \label{eq:operator-polynomial-bound}
\end{equation}
Consequently, $\overline{W(A)}$ is a $2$-spectral set for $A$.
\end{corollary}

\begin{proof}
The assertion is immediate if $p=0$.  Otherwise, write $d=\deg p$.
For a unit vector $x\in\mathcal H$, let
\[
  \mathcal M_x=\operatorname{span}\{x,Ax,\ldots,A^d x\},
  \qquad
  A_x=\left.\Pi_x A\right|_{\mathcal M_x},
\]
where $\Pi_x$ is the orthogonal projection onto
$\mathcal M_x$.  Induction gives
\begin{equation}
  A_x^k x=A^k x \qquad(0\leq k\leq d),
  \label{eq:Krylov-agreement}
\end{equation}
because $A^{k+1}x\in\mathcal M_x$ whenever $k<d$.  Moreover,
$W(A_x)\subset W(A)$.  Theorem~\ref{thm:main}, applied on the
finite-dimensional space $\mathcal M_x$, therefore gives
\[
  \lVert p(A)x\rVert_{\mathcal H}
  =\lVert p(A_x)x\rVert_{\mathcal H}
  \leq\lVert p(A_x)\rVert
  \leq2\sup_{z\in W(A)}\lvert p(z)\rvert.
\]
Taking the supremum over unit vectors $x$ proves~%
\eqref{eq:operator-polynomial-bound}.

Put $K=\overline{W(A)}$.  The standard inclusion
$\sigma(A)\subset K$ shows that $r(A)$ is defined for every rational
function $r$ pole-free on $K$.  The polynomial-approximation
argument in the proof of Theorem~\ref{thm:main}, now applied to $K$,
gives polynomials $p_j$ such that $p_j(A)\to r(A)$ and
\[
  \sup_{z\in W(A)}\lvert p_j(z)\rvert
  \longrightarrow
  \sup_{z\in W(A)}\lvert r(z)\rvert.
\]
Passing to the limit in~\eqref{eq:operator-polynomial-bound} proves the
rational estimate and hence the spectral-set assertion.
\end{proof}

The conclusion also controls polynomial and rational approximation of
matrix functions.  Let $A\in\Mn$, let $f$ be holomorphic on a
neighborhood of $W(A)$, and let $s$ be either a polynomial or a
rational function pole-free on $W(A)$.  Then $f-s$ is holomorphic on
a neighborhood of $W(A)$, and Corollary~\ref{cor:holomorphic-bound}
gives
\begin{equation}
  \lVert f(A)-s(A)\rVert
  \leq2\max_{z\in W(A)}\lvert f(z)-s(z)\rvert.
  \label{eq:matrix-function-error}
\end{equation}
This is the sharp universal form of the numerical-range estimate that
underlies applications to GMRES and rational Krylov methods; see~%
\cite{ChoiGreenbaum2015,CrouzeixGreenbaum2019,ChenGreenbaumTrogdon2025}
for the role of spectral sets in those settings.

\section{Discussion and further directions}
\label{sec:outlook}

Theorem~\ref{thm:main} proves the scalar-valued Crouzeix
conjecture.  The most immediate strengthening is the uniform
matrix-polynomial question: for every
$m\geq1$ and every matrix polynomial
$\mathcal F=[f_{ij}]\in\mathbb M_m(\C[z])$, does one have
\begin{equation}
  \lVert\mathcal F(A)\rVert
  \leq2\max_{z\in W(A)}\lVert\mathcal F(z)\rVert,
  \qquad \mathcal F(A)=[f_{ij}(A)].
  \label{eq:uniform-matrix-polynomial-extension}
\end{equation}
Here $\lVert\mathcal F(A)\rVert$ is the operator norm on $(\C^n)^m$,
whereas $\lVert\mathcal F(z)\rVert$ is the operator norm on $\C^m$.
The case $m=1$ is precisely the scalar theorem proved above;
uniformity in $m$ is an additional operator-algebraic requirement.  The
foundational matrix-level disk estimate is due to Okubo and Ando~%
\cite{OkuboAndo1975}.  Uniform matrix-level spectral-set bounds and
their relation to numerical radius are studied in~%
\cite{DavidsonPaulsenWoerdeman2018,HartzMcCarthy2026}; the constant
$1+\sqrt2$ at every matrix level follows from~%
\cite{CrouzeixPalencia2017}, while sharp estimates for
$\rho$-contractions with $1\leq\rho\leq2$ are developed in~%
\cite{SchwenningerDeVriesRho2025}.  The extension is
known with constant $2$ for weighted shift matrices~%
\cite{CrouzeixGreenbaum2025}.  For block-valued inputs, the coefficients
in the completion may occupy a larger operator algebra, so the
cancellation in~\eqref{eq:theta-contribution} would require a genuinely
operator-algebraic formulation.

The positive-real completion principle is independent of numerical
ranges: it starts with a Carath\'eodory function whose resolvent defect
lies in the adjoint algebra of an auxiliary simple-spectrum matrix
sharing an eigenbasis with the matrix in the resolvent.  A direct
infinite-dimensional extension of Theorem~\ref{thm:completion},
based for example on spectral measures or Riesz bases, could reveal
equality cases or uniform matrix-level functional-calculus structure.
The main challenge is a stable functional-calculus formulation of the
auxiliary-basis correction, valid beyond finite diagonal spectral
models.
Related questions for unbounded numerical ranges are studied in~%
\cite{CrouzeixUnbounded2025}, and related abstract spectral-constant
problems in~\cite{SchwenningerDeVries2024,BlazhkoEtAl2025}.

Double-layer methods beyond numerical ranges were developed in~%
\cite{CrouzeixGreenbaum2019}, while configuration-constant refinements
were introduced in~\cite{MalmanEtAl2025}; see also~%
\cite{SchwenningerDeVries2025} for a comparison of double-layer
approaches.
Annular double-layer kernels are studied in~\cite{JuryTsikalas2024},
while~\cite{AglerLykovaYoung2025} develops a complementary annular
realization and norm theory; elliptic domains admit structured dilation
models~\cite{AglerLykovaYoung2024}.  Whenever a companion transform
lands in a controlled adjoint algebra, the kernel-sampling argument may
replace a triangle-inequality estimate and thereby offer a route to
sharp constants for other geometric spectral sets.

With the notation of the proof of Theorem~\ref{thm:completion},
if $\lVert\widetilde T\rVert=2$ and $e\in\C^n$ satisfies
$\lVert e\rVert_2=1$ and $\lVert\widetilde Te\rVert_2=2$, then equality
throughout the last
operator-inequality chain in the proof of
Theorem~\ref{thm:completion} forces
\[
  \widehat Pe=2e,
  \qquad
  \widehat P\widetilde Te=\widetilde Te.
\]
Analyzing these vector saturation conditions may classify extremal
pairs $(A,p)$ and quantify stability near the nilpotent example~%
\eqref{eq:sharp-example}.  Such a
classification must allow the nonuniqueness phenomena found in~%
\cite{Li2020}.

Finally, O'Loughlin and Rani extend the Crouzeix--Palencia framework to
scaled $q$-numerical ranges, $q\in[0,1]$, and propose associated
spectral-set problems~\cite{OLoughlinRani2026}.  These provide a natural
testing ground for investigating a constant-$2$ analogue.

\section*{AI-assistance disclosure}

OpenAI ChatGPT contributed the idea of sampling the matrix Herglotz
kernel at half the conjugate eigenvalues and adding the origin sample
that cancels the adjoint-algebra correction; this appears in the proof
of Theorem~\ref{thm:completion}, especially~%
\eqref{eq:eigenvalue-samples}--\eqref{eq:theta-contribution}.  It also
assisted with exposition, bibliography, and typesetting.  The human
author assumes full responsibility for the correctness of all
mathematical statements and for the integrity and accuracy of all
citations.

\section*{Acknowledgments}

The author thanks Elijah Winners for suggesting the auxiliary-basis
simplification used in the proof of
Theorem~\ref{thm:completion}.

\reftitle{References}
\bibliography{the_numerical_range_is_a_2_spectral_set}

\end{document}
